\subsection{Préliminaires I, 3--5 : restrictions et types des morphismes}
\label{note:base0-sheaves}
Une section en 3.1.6 est prise sur $V$ ; la base de 3.2.2 est une base
de la topologie de l'ouvert $U$. Pour le morphisme d'espaces annelés de
4.1.1, l'homomorphisme structural porte sur les faisceaux d'anneaux
$\mathcal B\to\mathcal A$. En 4.4.8, le second morphisme de modules
est $v'$. Ces domaines déterminent les compositions considérées.

La présentation locale de 5.1.3 se déroule sur un seul ouvert $U$ :
le conoyau et les deux modules libres sont tous restreints à $U$.
En 5.2.4, l'exactitude à droite est celle du foncteur d'image réciproque
$f_U^*$, qui comporte un produit tensoriel. Le voisinage de 5.3.9
est un voisinage du point $x$.
En 5.4.5, $\mathcal L$ désigne le module localement libre ; en 5.4.7,
l'expression dans une trivialisation sur $U$ concerne les germes $x\in U$.
Dans 5.5.4, le faisceau est $\mathcal F$, et les zéros de la section sur
$U$ sont des points de $U$. Les bases et voisinages de 5.5.5 sont
également considérés dans cet ouvert $U$.

\subsection{Préliminaires I, 5.4 : modules inversibles}
\label{note:base0-invertible}
L'homomorphisme par homothéties
$\mathcal O_X\to\mathcal H\!om(\mathcal F,\mathcal F)$ existe pour
tout module $\mathcal F$. Il n'est pas généralement un isomorphisme :
pour $\mathcal F=\mathcal O_X^2$, les endomorphismes sont des matrices,
et les homothéties sont seulement les matrices scalaires.
Pour un module inversible $\mathcal L$, on se ramène localement à
$\mathcal L=\mathcal O_X$, et l'homomorphisme devient un isomorphisme.
C'est cette distinction qu'utilise 5.4.3.

Les changements de trivialisation d'un module inversible sont des
unités. Le groupe de cocycles de 5.4.7 est donc
$H^1(\mathfrak U,\mathcal O_X^*)$, avec le groupe multiplicatif des
unités, et non avec le faisceau additif $\mathcal O_X$.
En 5.4.8, le foncteur $f^*$ est restreint à la catégorie des modules
inversibles ; il définit ainsi l'application sur leurs classes
d'isomorphisme.

\subsection{Préliminaires I, 7 : gradués et complétions}
\label{note:base0-completion}
Dans 7.2.9, les générateurs initiaux sont considérés modulo le premier
terme positif $M^{(1)}$ de la filtration. Avec la convention
$M_n=M/\mathfrak J^{n+1}M$, le noyau correspondant à
$\mathfrak J^nM$ est celui de $M\to M_{n-1}$.
L'indice de la limite projective de 7.3.1 est $n$.
En 7.5.1, la composante $y_\mu$ appartient à $B_\mu$ ; le morphisme
de coefficients qui lui est appliqué est donc $\varphi_{\mu,\alpha}$.
La notation $\mathfrak J_\lambda$ au but conserve sa convention
d'image dans le quotient considéré.

Dans 7.5.5, les $c_i$ relèvent des générateurs du quotient, tandis
que les $a_j$ engendrent le noyau de carré nul. Leur réunion engendre
l'algèbre : on exprime d'abord l'image d'un élément dans le quotient,
puis la différence dans le noyau. Les seuls produits $c_ia_j$
resteraient dans le noyau et ne pourraient engendrer le quotient.
La propriété universelle de 7.7.6 comporte bien existence et unicité
de l'homomorphisme continu vers l'algèbre séparée complète : existence
par prolongement, unicité par densité.
\par
